% Code-aligned UPB extension. The LPB method and its shared implementation
% details are not repeated here.

\section{Extension to LLM utility evaluation and upper predictive bounds}
\label{sec:utility_upb}

Our framework also constructs upper predictive bounds (UPBs) on an event
time. One important application is LLM utility evaluation: an adversarial
attacker can be replaced by a helper agent that guides the target model toward
successful task completion. For example, in AutoIF~\citep{dong2025selfplay},
the helper iteratively assists the target model in satisfying complex
instructions. In this application, $Y_i(t)=1$ denotes successful task
completion, $T_i$ is the corresponding time-to-success, and a smaller UPB
certifies that success should occur within fewer interactions. The UPB
construction itself is generic and can be applied to any event represented by
$T_i$.

The observable interaction horizon is capped at $\maxl$. We encode ``no event
by turn $\maxl$'' by the additional value $\maxl+1$, which is treated as
infinity, and define
\[
    T_i^{(\maxl)}:=\min\{T_i,\maxl+1\}.
\]
The target guarantee is
\begin{equation}
    \mathbb P\!\left\{
        T_{\mathrm{test}}^{(\maxl)}
        \leq \widehat U(X_{\mathrm{test}})
        \ \middle|\ \Dall
    \right\}
    \geq 1-\alpha.
    \label{eq:upb_target}
\end{equation}
Thus, $\maxl$ remains a finite bound that can fail to cover a trajectory with
no event through turn $\maxl$, whereas $\maxl+1$ is always valid. Acquisition
is nevertheless truncated at $\maxl$; no interaction is acquired at time
$\maxl+1$.

\paragraph{Candidate-specific UPB errors.}
For a candidate $f_i(\tau):=\fhat_\tau(X_i)$, define
\begin{equation}
    A_i(\tau)
    :=\mathbb I\{f_i(\tau)<\maxl+1\}
      \mathbb I\{T_i>f_i(\tau)\}.
    \label{eq:upb_target_indicator}
\end{equation}
Unlike LPB miscoverage, this error is resolved when either the event is
observed or the acquisition reaches the candidate bound. For a finite
candidate, define its resolution time and resolution propensity by
\begin{equation}
    r_i(\tau):=\min\{T_i,f_i(\tau)\},
    \qquad
    \rho_i(t):=\prod_{s=1}^{t}P_i(s),
    \qquad
    \pi_i(\tau):=\rho_i(r_i(\tau)),
    \label{eq:upb_candidate_weight}
\end{equation}
with $\rho_i(0):=1$. Here $P_i(t)$ is the conditional probability of
acquiring turn $t$ given that the preceding prefix has been reached. If
$Z_i(t)$ is the corresponding Bernoulli continuation decision, the nested
reach indicator is
\begin{equation}
    R_i(0):=1,
    \qquad
    R_i(t):=\prod_{s=1}^{t}Z_i(s)
           =\mathbb I\{\min(C_i,T_i)\geq t\}.
    \label{eq:upb_reach_indicator}
\end{equation}
Consequently,
$\mathbb E_{\rm acq}[R_i(t)\mid X_i,H_i,T_i]=\rho_i(t)$ once the policy is
frozen. On the fully observed Phase-I rows, $P_i(t)=R_i(t)=1$ for every
relevant turn. The ordinary-HT contribution for a finite candidate is
$R_i(r_i(\tau))A_i(\tau)/\pi_i(\tau)$; the $\maxl+1$ candidate has
deterministic contribution zero.

\paragraph{Sequential augmented miscoverage estimator.}
Our implemented UPB calibrator uses the survival-model output at every reached
prefix. For a finite candidate $f=f_i(\tau)$, let
\begin{equation}
    m_{i,0}(f):=\widehat{\mathbb P}(T_i>f\mid H_i(0)),
    \qquad
    m_{i,t}(f):=\widehat{\mathbb P}(T_i>f\mid H_i(t))
    \quad (t<r_i(\tau)),
\end{equation}
and set $m_{i,r_i(\tau)}(f):=A_i(\tau)$. The prediction path is a frozen,
predictable function of the observed histories; it need not be correctly
specified. The row-wise sequential augmented Horvitz--Thompson (AHT)
contribution is
\begin{equation}
    \widetilde A_i(\tau)
    :=m_{i,0}(f_i(\tau))
      +\sum_{t=1}^{r_i(\tau)}
       \frac{R_i(t)}{\rho_i(t)}
       \left[m_{i,t}(f_i(\tau))-m_{i,t-1}(f_i(\tau))\right].
    \label{eq:upb_sequential_aht}
\end{equation}
For $f_i(\tau)=\maxl+1$, we instead set
$\widetilde A_i(\tau):=0$.

Equation~\eqref{eq:upb_sequential_aht} starts from the model's initial
prediction and adds an inverse-probability correction for every prediction
update revealed by acquisition. The corrections telescope in expectation to
the resolved error $A_i(\tau)$, so misspecification affects efficiency but not
design validity. This construction is related to prediction-powered
inference~\citep{angelopoulos2023prediction} and augmented
Horvitz--Thompson estimation~\citep{Horvitz01121952,robins1994estimation,wang2001augmented}.
If
\begin{equation}
    m_{i,t}(f)=m_{i,0}(f)
    \quad\text{for all }t<r_i(\tau),
    \label{eq:upb_no_intermediate_updates}
\end{equation}
then only the terminal update is nonzero and
\begin{equation}
    \widetilde A_i(\tau)
    =m_{i,0}(f)
     +\frac{R_i(r_i(\tau))}{\pi_i(\tau)}
       \{A_i(\tau)-m_{i,0}(f)\},
    \label{eq:upb_terminal_residual_aht}
\end{equation}
which is the terminal residual-AHT estimator.

\begin{proposition}[Validity of the UPB acquisition estimator]
\label{prop:valid_upb_weights}
Fix a candidate $\tau$ and condition on the data used to freeze the
acquisition policy and on the complete trajectory $(X_i,H_i,T_i)$. If the
relevant reach probabilities are positive, then
\begin{equation}
    \mathbb E_{\rm acq}\!\left[\widetilde A_i(\tau)\right]
    =A_i(\tau).
    \label{eq:valid_upb_weight_identity}
\end{equation}
This identity holds for any predictable sequence $m_{i,t}$. Moreover, for a
finite candidate,
\begin{equation}
    \operatorname{Var}_{\rm acq}\!\left(\widetilde A_i(\tau)\right)
    =\sum_{t=1}^{r_i(\tau)}
      \left(\frac{1}{\rho_i(t)}-\frac{1}{\rho_i(t-1)}\right)
      \left(A_i(\tau)-m_{i,t-1}(f_i(\tau))\right)^2.
    \label{eq:upb_sequential_aht_variance}
\end{equation}
\end{proposition}

\begin{proof}
Suppress $i$ and $\tau$, write $r=r_i(\tau)$, $A=A_i(\tau)$,
$m_t=m_{i,t}(f_i(\tau))$, $Q_t=R_i(t)/\rho_i(t)$, and $Q_0=1$.
For the acquisition filtration $\{\mathcal G_t\}$,
\[
    \mathbb E_{\rm acq}[Q_t\mid\mathcal G_{t-1}]=Q_{t-1},
    \qquad
    \mathbb E_{\rm acq}[Q_t]=1.
\]
Using $m_r=A$ in~\eqref{eq:upb_sequential_aht} therefore gives
\[
    \mathbb E_{\rm acq}[\widetilde A_i(\tau)]
    =m_0+\sum_{t=1}^{r}(m_t-m_{t-1})=A.
\]
For the variance, define $D_t:=Q_t-Q_{t-1}$. Summation by parts gives
\[
    \widetilde A_i(\tau)-A
    =\sum_{t=1}^{r}D_t(A-m_{t-1}).
\]
The $D_t$ are orthogonal martingale differences and, using
$R_i(t)R_i(t-1)=R_i(t)$,
\[
    \mathbb E_{\rm acq}[D_t^2]
    =\frac{1}{\rho_i(t)}-\frac{1}{\rho_i(t-1)}.
\]
Substitution proves~\eqref{eq:upb_sequential_aht_variance}. The
$\maxl+1$ case is immediate because both $A_i(\tau)$ and
$\widetilde A_i(\tau)$ are zero.
\end{proof}

We estimate the candidate miscoverage rate by
\begin{equation}
    \widehat\alpha_{\rm UPB}(\tau)
    :=\frac{1}{|\calI|}\sum_{i\in\calI}\widetilde A_i(\tau),
    \qquad
    \widehat c_{\rm UPB}(\tau):=1-\widehat\alpha_{\rm UPB}(\tau).
    \label{eq:upb_miscoverage_estimator}
\end{equation}
Let $\mathcal T_{\rm prior}=\{\tau_1<\cdots<\tau_J\}$ denote the finite
candidate grid no larger than $\tauprior$. Since empirical AHT estimates need
not be monotone, the implementation forms the suffix lower envelope of
coverage and selects
\begin{equation}
    \widehat\tau
    :=\min\left\{
        \tau_j\in\mathcal T_{\rm prior}:
        \min_{k\geq j}\widehat c_{\rm UPB}(\tau_k)\geq1-\alpha
    \right\}.
    \label{eq:upb_calibrated_tau}
\end{equation}
Equivalently, it requires
$\max_{k\geq j}\widehat\alpha_{\rm UPB}(\tau_k)\leq\alpha$.
The calibrated bound is $\widehat U(x):=\fhat_{\widehat\tau}(x)$. If no
finite candidate through $\tauprior$ satisfies the rule, the implementation
returns the deterministic fallback $\widehat U(x)=\maxl+1$.

\begin{theorem}[UPB coverage validity]
\label{thm:upb_coverage}
Suppose the sampling and independence conditions of
Theorem~\ref{thm:general_validity} hold for the UPB contributions above.
Assume $m_{i,t}(f)\in[0,1]$, $\rho_i(t)\geq\varepsilon>0$ at every relevant
prefix, and, uniformly over $\tau\in\mathcal T_{\rm prior}$,
\begin{equation}
    \mathbb E\!\left[\pi_i(\tau)^{-1}\right]\leq\bar w.
\end{equation}
Let $n:=|\calI|$ and define
\begin{equation}
    \Delta_{\rm UPB}
    :=\frac{\log(1/\delta)}{3n\varepsilon}
      +\sqrt{
        \frac{\log^2(1/\delta)}{9n^2\varepsilon^2}
        +\frac{2(\bar w-\alpha^2)\log(1/\delta)}{n}}
    .
    \label{eq:upb_coverage_slack}
\end{equation}
Then, with probability at least $1-\delta$ over the calibration sample and
acquisition randomization,
\begin{equation}
    \mathbb P\!\left\{
      T_{\rm test}^{(\maxl)}\leq\widehat U(X_{\rm test})
      \ \middle|\ \Dall
    \right\}
    \geq1-\alpha-\Delta_{\rm UPB}.
    \label{eq:upb_coverage_result}
\end{equation}
\end{theorem}

\begin{proof}
For any finite candidate, Proposition~\ref{prop:valid_upb_weights} gives
design unbiasedness. If $k$ is the last reached prefix, summation by parts
gives
\[
    \widetilde A_i
    =\frac{m_{i,k}}{\rho_i(k)}
      +\sum_{t=0}^{k-1}m_{i,t}
        \left(\frac{1}{\rho_i(t)}-\frac{1}{\rho_i(t+1)}\right),
\]
and hence $1-\varepsilon^{-1}\leq\widetilde A_i\leq\varepsilon^{-1}$.
Equation~\eqref{eq:upb_sequential_aht_variance} and
$|A_i-m_{i,t-1}|\leq1$ imply
\[
    \mathbb E[\widetilde A_i^2]
    =\mathbb E[\operatorname{Var}_{\rm acq}(\widetilde A_i)+A_i^2]
    \leq\mathbb E[\pi_i^{-1}]
    \leq\bar w.
\]
Applying the one-sided Bernstein argument from
Theorem~\ref{thm:general_validity}, with upper centered range
$\varepsilon^{-1}$, yields~\eqref{eq:upb_coverage_slack}. The suffix-envelope
selector in~\eqref{eq:upb_calibrated_tau} supplies the UPB analogue of the
monotone candidate restriction used in that theorem. The $\maxl+1$ fallback
has zero miscoverage by construction.
\end{proof}

\paragraph{UPB specialization of \ttmethod.}
For LPBs, Phase I uses the fixed reference level $\alpha$ because the final
calibrated level is unknown during acquisition. For UPBs, the implementation
instead selects a reference level from the initial, label-free survival-model
predictions. Define
\begin{equation}
    \widehat a_0(\tau)
    :=\frac{1}{|\calI|}\sum_{i\in\calI}
      \mathbb I\{f_i(\tau)\leq\maxl\}
      \widehat{\mathbb P}(T_i>f_i(\tau)\mid H_i(0)).
\end{equation}
The implementation chooses
\begin{equation}
    \tau_{\rm ref}
    :=\min\left\{
      \tau_j\in\mathcal T_{\rm prior}:
      \min_{k\geq j}\{1-\widehat a_0(\tau_k)\}\geq1-\alpha
    \right\},
    \label{eq:upb_reference_level}
\end{equation}
using $\max\mathcal T_{\rm prior}$ if the set is empty, and then freezes
\begin{equation}
    f_i^\star:=\fhat_{\tau_{\rm ref}}(X_i).
\end{equation}
This selection uses the initial model outputs for all calibration trajectories
but no trajectory labels or acquired interactions. It respects the reversed
direction of UPB calibration and adapts to the requested target coverage.

The deployed score remains the causal current-event probability
$S_i(t)=\widehat h_i(t)$ from~\eqref{eq:lpb_impl_hazard_score}; only the
task-specific objective coefficients change. With the same global
regularization $\delta$ as in the LPB implementation, the UPB coefficients are
\begin{equation}
\begin{split}
    a_i^{\rm UPB}(t)
    :=\frac{1}{1+\delta}\biggl[&
      \mathbb I\{t=f_i^\star,\ f_i^\star\leq b_i,\ f_i^\star\leq\maxl\}
      \bigl(1-\widehat h_i(t)\bigr)\\
      &+\delta\,\mathbb I\{t\leq b_i\}\widehat h_i(t)
    \biggr].
    \label{eq:upb_dapro_soft_coefficients}
\end{split}
\end{equation}
At a finite reference bound $f_i^\star$, the factor
$1-\widehat h_i(f_i^\star)$ is the estimated conditional probability of
surviving beyond that bound given that its pre-decision prefix has been
reached. The value $\maxl+1$ has deterministic zero miscoverage and therefore
contributes no target-specific mass.

The implemented UPB DAPRO policy substitutes
$a_i^{\rm UPB}(t)$ for the LPB coefficients in the same inverse-reach
optimization:
\begin{equation}
\begin{aligned}
    \min_M\quad &
    \frac{1}{N_1}\sum_{i\in\calIone}\sum_{t=1}^{b_i}
    \frac{a_i^{\rm UPB}(t)}{\prod_{s=1}^{t}M_s(S_i(s))}\\
    \text{s.t.}\quad &
    \frac{1}{N_1}\sum_{i\in\calIone}\bfunc(M(S_i))\leq\Bbar,\\
    &S_i(t)\leq S_j(t)
      \Longrightarrow M_t(S_i(t))\leq M_t(S_j(t))
      \quad\text{for every active }i,j,t.
    \label{eq:upb_dapro_objective}
\end{aligned}
\end{equation}
This is the endpoint soft-mass allocation proxy used by the current code. It
is an ordinary-HT-style inverse-reach proxy; it is \emph{not} the exact
sequential-AHT variance in~\eqref{eq:upb_sequential_aht_variance}. Sequential
AHT is used downstream to exploit every reached prediction update while the
allocation objective remains fixed. Thus, the estimator comparison can change
the estimator without silently changing the acquired trajectories.

All remaining components---the two-bin score-to-probability map, numerical
optimization, cumulative-reach budget correction, positivity floor, and the
optional independent CRC budget controller---are shared with the LPB
implementation in Section~\ref{sec:proposed_alg}. The reported UPB DAPRO
policy is history-adaptive: it recomputes $S_i(t)$ and makes a new continuation
decision after each observed interaction.
